\documentclass{article}%
\usepackage[T1]{fontenc}%
\usepackage[utf8]{inputenc}%
\usepackage{lmodern}%
\usepackage{textcomp}%
\usepackage{lastpage}%
\usepackage{geometry}%
\geometry{margin=1in,headheight=14pt,headsep=25pt}%
\usepackage{hyperref}%
\usepackage{enumitem}%
\usepackage{fancyhdr}%
\usepackage{xcolor}%
\usepackage{url}%
\usepackage{breakurl}%
\usepackage{amsmath}%
\usepackage{amssymb}%
\usepackage{mathtools}%
\usepackage{amsthm}%
\usepackage{thmtools}%
%

            \hypersetup{
                colorlinks=true,
                linkcolor=blue,
                filecolor=magenta,
                urlcolor=blue,
                breaklinks=true,
                bookmarks=true,
                pdfborder={0 0 0}
            }
            \urlstyle{same}
            \Urlmuskip=0mu plus 1mu  % URL breaking in nicer spots
            
            % Math configuration
            \everymath{\displaystyle}
            \setlength{\jot}{10pt}  % Increase spacing between equations
            
            \pagestyle{fancy}
            \fancyhf{}
            \rhead{Generated on \today}
            \cfoot{\thepage}
            
            % Better Unicode support
            \usepackage[utf8]{inputenc}
            \usepackage{newunicodechar}
            \usepackage[breaklinks=true]{hyperref}
            \setlength{\itemsep}{1em} % Adds spacing between items
            
            % Configure itemize settings globally
            \setlist[itemize]{nosep}
            \setlist[itemize]{leftmargin=*}
            
            % Remove section numbering
            \renewcommand{\thesection}{}
            \renewcommand{\thesubsection}{}
        %
\lhead{Lecture: 2024{-}09{-}05{-}SimplexEfficiency}%
\title{2024{-}09{-}05{-}SimplexEfficiency}%
\author{Generated by Scribe.AI}%
\date{\today}%
%
\begin{document}%
\normalsize%
\maketitle%
\section{Slide 1}%
\label{sec:Slide1}%
\subsection{Content}%

\underline{\text{Efficiency of Simplex Method}} \\
\text{Worst case vs Average Case} \\
\bullet \text{ (Dantzig) } m < 50, \quad m+n < 200 \\
\#(\text{iterations}) \le \frac{3m}{2} \text{ (usually)}, \quad \text{rarely } \approx 3m
%
\subsection{Description}%
This slide discusses the efficiency of the simplex method in the context of linear programming. The title "Efficiency of Simplex Method" is underlined.  It compares the worst-case and average-case performance.  According to Dantzig, for problems with $m < 50$ and $m + n < 200$, the number of iterations is usually less than or equal to $\frac{3m}{2}$, but rarely approaches $3m$.  Two tables are included, likely from a source like Dantzig's work, showing empirical data on the average number of iterations required by two different simplex rules (largest-coefficient and largest-increase) for various problem sizes ($m$ and $n$).  The tables provide evidence supporting the claim about the average-case performance of the simplex method. The context of the course (Linear Programming and Optimization Techniques) is crucial for understanding the significance of this analysis, as it addresses the practical efficiency of a fundamental linear programming algorithm.%
\subsection{Figures}%
\begin{figure}[h]%
\centering%
\begin{minipage}{\textwidth}%
\centering%
\begin{minipage}{0.15\textwidth}%
\centering%
\includegraphics[width=0.8\textwidth]{/Users/ashoksaravanan/Coding/ScribeLec/server/output/MA 421/lectures/2024-09-05-SimplexEfficiency/figures/1.1.png}%
\caption{Two tables showing average numbers of iterations required by the largest-coefficient rule and the largest-increase rule for different values of $m$ and $n$.}%
\end{minipage}%
\end{minipage}%
\end{figure}

%
\newpage%
\section{Slide 2}%
\label{sec:Slide2}%
\subsection{Content}%

\underline{\text{Worst Case Example: Klee-Minty (1970)}} \\
\text{The example is quite simple to state:} \\
\text{maximize } \sum_{j=1}^n 2^{n-j} x_j \\
\text{(4.1) subject to } \sum_{j=1}^i 2^{i-j} x_j + x_{i+1} \le 100^{i-1}, \quad i = 1, 2, \dots, n \\
x_j \ge 0 \quad j = 1, 2, \dots, n \\
\text{It is instructive to look closely at the constraints. The first three constraints are:} \\
x_1 \le 1 \\
4x_1 + x_2 \le 100 \\
8x_1 + 4x_2 + x_3 \le 10,000
%
\subsection{Description}%
This slide presents a worst-case example for the simplex method, known as the Klee-Minty cube. The title "Worst Case Example: Klee-Minty (1970)" is underlined. The example is a linear program in the form of a maximization problem with $n$ variables and $n$ constraints. The objective function is $\sum_{j=1}^n 2^{n-j} x_j$. The constraints are given by $\sum_{j=1}^i 2^{i-j} x_j + x_{i+1} \le 100^{i-1}$ for $i = 1, 2, \dots, n$, and $x_j \ge 0$ for $j = 1, 2, \dots, n$. The slide then highlights the first three constraints of this example, showing how they are structured.  The context of the course (Linear Programming and Optimization Techniques) indicates that this example is used to illustrate that the simplex method can have exponential time complexity in the worst case. The Klee-Minty cube is a classic example demonstrating that the simplex method, while efficient in practice, can take an exponential number of iterations in certain pathological cases. The slide likely serves as an introduction to the analysis of the simplex method's worst-case complexity.%
\subsection{Figures}%
\begin{figure}[h]%
\centering%
\begin{minipage}{\textwidth}%
\centering%
\begin{minipage}{0.15\textwidth}%
\centering%
\includegraphics[width=0.8\textwidth]{/Users/ashoksaravanan/Coding/ScribeLec/server/output/MA 421/lectures/2024-09-05-SimplexEfficiency/figures/2.1.png}%
\caption{No figure present.}%
\end{minipage}%
\end{minipage}%
\end{figure}

%
\newpage%
\section{Slide 3}%
\label{sec:Slide3}%
\subsection{Content}%

\underline{\text{Worst Case Example: Klee-Minty (1970)}} \\
\text{[v]} \\
\text{maximize } \sum_{j=1}^n 2^{n-j} x_j - \frac{1}{2} \sum_{j=1}^n 2^{n-j} \beta_j \\
\text{(4.2) subject to } 2 \sum_{j=1}^{i-1} 2^{i-1-j} x_j + x_i \le \sum_{j=1}^{i-1} 2^{i-1-j} \beta_j + \beta_i, \quad i = 1, 2, \dots, n \\
x_j \ge 0 \quad j = 1, 2, \dots, n \\
1 \ll \beta_1 \ll \beta_2 \ll \dots \ll \beta_n \\
\underline{\text{Largest Coefficient Rule: require } 2^n - 1 \text{ steps}} \\
\underline{\text{Smallest Coefficient Rule: just } n \text{ steps}}
%
\subsection{Description}%
This slide presents the Klee-Minty example (1970), a worst-case scenario for the simplex method in linear programming. The title "Worst Case Example: Klee-Minty (1970)" is underlined.  The example is a linear program formulated as a maximization problem with $n$ variables ($x_j$) and $n$ constraints. The objective function is a weighted sum of the variables $x_j$, subtracted by a weighted sum of parameters $\beta_j$. The constraints are inequalities involving weighted sums of the variables $x_j$ and parameters $\beta_j$. The parameters $\beta_j$ satisfy the condition $1 \ll \beta_1 \ll \beta_2 \ll \dots \ll \beta_n$, indicating a significant difference in magnitude between consecutive $\beta_j$ values.  The slide then highlights the performance of two different pivot rules: the Largest Coefficient Rule, which requires $2^n - 1$ steps (exponential complexity), and the Smallest Coefficient Rule, which requires only $n$ steps (linear complexity).  The underlining emphasizes the contrasting performance of these rules. The context of the course (Linear Programming and Optimization Techniques) is crucial because this example demonstrates that the simplex method's efficiency is highly dependent on the chosen pivot rule and can exhibit exponential time complexity in the worst case, despite its generally good performance in practice.%
\subsection{Figures}%
\begin{figure}[h]%
\centering%
\begin{minipage}{\textwidth}%
\centering%
\begin{minipage}{0.15\textwidth}%
\centering%
\includegraphics[width=0.8\textwidth]{/Users/ashoksaravanan/Coding/ScribeLec/server/output/MA 421/lectures/2024-09-05-SimplexEfficiency/figures/3.1.png}%
\caption{No figure present.}%
\end{minipage}%
\end{minipage}%
\end{figure}

%
\newpage%
\end{document}